Similar agent reports can restate one observation, describe independent executions,
or contain complementary requirements whose differences determine a decision.
We develop a controlled benchmark for studying these distinctions in hierarchical
language-model systems. The primary task is to select the feasible candidate
schedule with minimum makespan from a supplied set. Executable verifiers determine
constraint satisfaction and the correct decision, and an independently implemented
checker validates small integer instances. Paired interventions change numerical
limits, scoped identities, or dependencies; a separate rule-workflow population
tests quantifiers, conditions, and explicit revisions. Frozen reports support
comparisons between single presentations, exact copies, and claim-preserving
paraphrases under flat and summary-only hierarchical aggregation. The protocol
compares ordinary synthesis, required-task checklists, and scope-aware consolidation
with common information access and output contracts. This working manuscript
currently reports development and interface validation rather than completed
held-out hypothesis tests. Initial execution exposed formatting failures and
insufficient baseline task competence. After interface validation, the completed
one-model pilot showed ceiling flat scheduling accuracy and beneficial repetition
effects in hierarchical scheduling, with mixed secondary-workflow findings.
We document these limitations, preserve unsuccessful outputs, and
specify paired cluster-level analyses and additional controls needed to distinguish
repetition effects from information loss, missing coverage, and resource changes.
Human scientific review and the remaining experimental stages are required before
submission.
